\paragrafos{Διχοτόμος γωνίας}{Β.3.3, «Διχοτόμος γωνίας»}

\begin{anagke}
Στο διάλειμμα ο Νίκος και η Αγγελική μοιράζονται ένα κομμάτι πίτσα σε σχήμα «φέτας» με μυτερή γωνία. «Θέλω ακριβώς τη μισή», λέει η Αγγελική, «και όχι με το μάτι». Ο Νίκος κοιτάζει τη φέτα: «Αν κόψω από την μύτη προς την άκρη, πώς θα ξέρω ότι τα δύο κομμάτια είναι ίσα;». Η Φανή το σημειώνει: ζητούμενο είναι μια ευθεία από την κορυφή που χωρίζει τη γωνία σε δύο ίσες γωνίες.

Η Ρήνη βρίσκει και ένα δεύτερο πρόβλημα στο διπλανό γήπεδο: «Το καφενείο θέλει να βάλει ένα τραπέζι σε μια γωνία, ώστε να απέχει ίσα από τους δύο τοίχους». Ο Ηλίας υποψιάζεται ότι οι δύο προβλήματα είναι ένα: η ευθεία που κόβει τη γωνία στη μέση είναι και η ευθεία των σημείων που ισαπέχουν από τις πλευρές της.

Σε αυτή την παράγραφο μαθαίνουμε τι είναι η διχοτόμος μιας γωνίας, ποια ιδιότητα έχουν τα σημεία της και πώς τη σχεδιάζουμε με μοιρογνωμόνιο, με δίπλωση και με διαβήτη.
\end{anagke}

\begin{lexi}{διχοτόμος, απόσταση}{διχοτόμος}
Η \lexiword{διχοτόμος} είναι αυτή που «κόβει δίχα», δηλαδή στα δύο (δίχα + τέμνω): χωρίζει τη γωνία σε δύο ίσα μέρη. Η \lexiword{απόσταση} σημείου από ευθεία είναι το μήκος του κάθετου τμήματος που φέρνουμε από το σημείο προς την ευθεία.
\end{lexi}

\begin{metafrasi}
Η ημιευθεία που κόβει τη γωνία στη μέση. & η διχοτόμος $\mathrm{Οδ}$ της $x\hat{\mathrm{Ο}}y$\\
Οι δύο μισές γωνίες είναι ίσες. & $x\hat{\mathrm{Ο}}\delta=\delta\hat{\mathrm{Ο}}y$\\
Κάθε μισή είναι το μισό της γωνίας. & $x\hat{\mathrm{Ο}}\delta=\dfrac{x\hat{\mathrm{Ο}}y}{2}$\\
Το σημείο $\mathrm{Α}$ έχει ίσες αποστάσεις από τις δύο πλευρές. & $\mathrm{ΑΒ}=\mathrm{ΑΓ}$\\
\end{metafrasi}

\tmima{Η θεωρία}

\begin{orismos}{διχοτόμος γωνίας}
\textbf{Διχοτόμος} μιας γωνίας $x\hat{\mathrm{Ο}}y$ λέγεται η ημιευθεία $\mathrm{Οδ}$ που έχει αρχή την κορυφή της γωνίας και τη χωρίζει σε δύο ίσες γωνίες: $x\hat{\mathrm{Ο}}\delta=\delta\hat{\mathrm{Ο}}y$.
\end{orismos}

Κάθε γωνία έχει μία και μόνο μία διχοτόμο. Κάθε μία από τις δύο μισές γωνίες έχει μέτρο ίσο με το μισό της αρχικής.

\begin{center}
\begin{tikzpicture}[scale=0.95]
\draw[thick] (5,0)--(0,0)--({5*cos(60)},{5*sin(60)});
\draw[gblue,thick] (0,0)--({5*cos(30)},{5*sin(30)});
\fill (0,0) circle (2pt) node[below left] {$\mathrm{Ο}$};
\node[right] at (5,0) {$x$}; \node[above] at ({5*cos(60)},{5*sin(60)}) {$y$}; \node[right,gblue] at ({5*cos(30)},{5*sin(30)}) {$\delta$};
\draw[gblue] (1.1,0) arc(0:30:1.1); \draw[gblue] (1.4,0) arc(0:30:1.4); \draw[gwarmD] (1.3,0.75) arc(30:60:1.5); \draw[gwarmD] (1.6,0.92) arc(30:60:1.85);
\end{tikzpicture}
\end{center}

Θυμόμαστε ότι \textbf{απόσταση} ενός σημείου $\mathrm{Α}$ από μια ευθεία $\varepsilon$ είναι το μήκος του κάθετου τμήματος $\mathrm{ΑΒ}$ που φέρνουμε από το $\mathrm{Α}$ προς την $\varepsilon$.

\begin{idiotita}{σημεία της διχοτόμου}
Κάθε σημείο της διχοτόμου μιας γωνίας \textbf{ισαπέχει} από τις πλευρές της.
\end{idiotita}

Ισχύει και το αντίστροφο: κάθε σημείο στο εσωτερικό της γωνίας που ισαπέχει από τις πλευρές της ανήκει στη διχοτόμο της. Έτσι η διχοτόμος είναι ακριβώς το σύνολο των σημείων που ισαπέχουν από τις πλευρές της γωνίας, όπως η μεσοκάθετος είναι το σύνολο των σημείων που ισαπέχουν από δύο σημεία.

\begin{center}
\begin{tikzpicture}[scale=0.95]
\draw[thick] (5,0)--(0,0)--({5*cos(60)},{5*sin(60)});
\draw[gblue,thick] (0,0)--({4.5*cos(30)},{4.5*sin(30)});
\fill (0,0) circle (2pt) node[below left] {$\mathrm{Ο}$};
\coordinate (P) at (2.944,1.7); \coordinate (B) at (2.944,0); \coordinate (G) at (1.472,2.55);
\fill (P) circle (2pt) node[above right] {$\mathrm{Α}$};
\draw[gwarmD,thick] (P)--(B) node[midway,right] {$\delta$}; \draw[gwarmD,thick] (P)--(G) node[midway,above right] {$\delta$};
\node[below] at (B) {$\mathrm{Β}$}; \node[left] at (G) {$\Gamma$};
\draw (2.744,0)--(2.744,0.2)--(2.944,0.2);
\draw (1.372,2.377)--(1.545,2.277)--(1.645,2.45);
\end{tikzpicture}
\end{center}

\tmimasmall{Πώς σχεδιάζουμε τη διχοτόμο}

\textbf{Με μοιρογνωμόνιο.} Μετράμε τη γωνία, διαιρούμε το μέτρο με το $2$ και σημειώνουμε αυτή τη γωνία πάνω στη μία πλευρά.

\textbf{Με δίπλωση.} Διπλώνουμε διαφανές χαρτί ώστε οι δύο πλευρές της γωνίας να συμπέσουν. Η πτυχή είναι η διχοτόμος.

\textbf{Με κανόνα και διαβήτη.} Με κέντρο την κορυφή $\mathrm{Ο}$ γράφουμε τόξο που τέμνει τις πλευρές στα $\mathrm{Κ}$ και $\Lambda$. Με κέντρα τα $\mathrm{Κ}$ και $\Lambda$ και την ίδια ακτίνα γράφουμε δύο τόξα που τέμνονται στο $\mathrm{Π}$. Η $\mathrm{ΟΠ}$ είναι η διχοτόμος.

\begin{center}
\begin{tikzpicture}[scale=0.95]
\draw[thick] (5,0)--(0,0)--({5*cos(70)},{5*sin(70)});
\fill (0,0) circle (2pt) node[below left] {$\mathrm{Ο}$};
\draw[gblue] (2,0) arc(0:70:2);
\coordinate (K) at (2,0); \coordinate (L) at ({2*cos(70)},{2*sin(70)});
\fill (K) circle (2pt) node[below] {$\mathrm{Κ}$}; \fill (L) circle (2pt) node[left] {$\Lambda$};
\draw[gwarmD] (K)++(55:1.8) arc(55:95:1.8);
\draw[gwarmD] (L)++(-25:1.8) arc(-25:15:1.8);
\fill (2.478,1.735) circle (2pt) node[above right] {$\mathrm{Π}$};
\draw[gblue,thick] (0,0)--(3.9,2.73);
\end{tikzpicture}
\end{center}

\begin{prosoxi}
\begin{itemize}[leftmargin=1.2em,itemsep=0.3ex,topsep=2pt]
\item Η διχοτόμος είναι ημιευθεία μέσα στη γωνία, με αρχή την κορυφή. Δεν είναι τμήμα που ενώνει δύο σημεία των πλευρών.
\item Η απόσταση σημείου από ευθεία μετριέται πάντα κατά το \emph{κάθετο} τμήμα, όχι κατά ένα λοξό.
\item Σημείο στο εσωτερικό της γωνίας με άνισες αποστάσεις από τις πλευρές δεν ανήκει στη διχοτόμο.
\item Οι δύο μισές γωνίες είναι ίσες σε \emph{μέτρο}, ακόμη κι αν στο σχέδιο οι πλευρές τους έχουν διαφορετικό μήκος.
\end{itemize}
\end{prosoxi}

\tmima{Λυμένα παραδείγματα}

\paradeigma{Η φέτα της πίτσας}
\tixero Η γωνία της κορυφής της φέτας είναι $60^\circ$.\par
\tithelo Με ποια γωνία ξεκινά η κοψιά που τη χωρίζει σε δύο ίσα μέρη;\par
\possindeo Η κοψιά είναι η διχοτόμος της γωνίας, άρα χωρίζει τη γωνία σε δύο ίσες γωνίες.\par
\apantisi $60^\circ:2=30^\circ$ με κάθε πλευρά.\par
\elegxos $30^\circ+30^\circ=60^\circ$.

\paradeigma{Η διχοτόμος και το μέτρο}
\tixero Η $\mathrm{Οδ}$ είναι διχοτόμος της $x\hat{\mathrm{Ο}}y$ και $x\hat{\mathrm{Ο}}\delta=28^\circ$.\par
\tithelo Το μέτρο της $x\hat{\mathrm{Ο}}y$.\par
\possindeo Οι δύο μισές γωνίες είναι ίσες: $\delta\hat{\mathrm{Ο}}y=x\hat{\mathrm{Ο}}\delta=28^\circ$.\par
\apantisi $x\hat{\mathrm{Ο}}y=28^\circ+28^\circ=56^\circ$.\par
\elegxos Η διχοτόμος των $56^\circ$ δίνει δύο γωνίες των $28^\circ$.

\paradeigma{Το ρολόι ως διχοτόμος}
\tixero Στις $4{:}00$ ο λεπτοδείκτης δείχνει το $12$ και ο ωροδείκτης το $4$. Κάθε ώρα είναι $30^\circ$.\par
\tithelo Ποιον αριθμό του ρολογιού δείχνει η διχοτόμος της γωνίας των δεικτών;\par
\possindeo Η γωνία των δεικτών είναι $4\cdot30^\circ=120^\circ$ και η μισή της είναι $60^\circ$, που αντιστοιχεί σε $60^\circ:30^\circ=2$ ώρες από το $12$.\par
\apantisi Η διχοτόμος δείχνει το $2$.\par
\elegxos Από το $12$ ως το $2$ είναι $60^\circ$ και από το $2$ ως το $4$ άλλες $60^\circ$.

\paradeigma{Το σημείο που ισαπέχει}
\tixero Μια γωνία $x\hat{\mathrm{Ο}}y=80^\circ$.\par
\tithelo Να βρω σημείο $\mathrm{Α}$ στο εσωτερικό της που να ισαπέχει από τις πλευρές της και να απέχει $5$ cm από την κορυφή.\par
\possindeo Τα σημεία που ισαπέχουν από τις πλευρές είναι τα σημεία της διχοτόμου. Φέρνω τη διχοτόμο (γωνία $40^\circ$ με κάθε πλευρά) και παίρνω πάνω της σημείο $\mathrm{Α}$ με $\mathrm{ΟΑ}=5$ cm.\par
\apantisi Το $\mathrm{Α}$ είναι πάνω στη διχοτόμο, σε απόσταση $5$ cm από το $\mathrm{Ο}$.\par
\elegxos Μετράω τις κάθετες από το $\mathrm{Α}$ προς τις πλευρές: βγαίνουν ίσες.

\tmima{Ασκήσεις}

\zoni{Το κατάλαβα}{διχοτόμος και ισαπόσταση}

\ex Βρες το μέτρο κάθε μισής γωνίας όταν η γωνία είναι: α) $70^\circ$ \quad β) $124^\circ$ \quad γ) $36^\circ$ \quad δ) $90^\circ$ \quad ε) $180^\circ$.

\ex Σχεδίασε μια γωνία $x\hat{\mathrm{Ο}}y=64^\circ$ και φέρε τη διχοτόμο της: α) με μοιρογνωμόνιο, β) με δίπλωση διαφανούς χαρτιού. Συμπίπτουν;

\ex Η $\mathrm{Οδ}$ είναι διχοτόμος της $x\hat{\mathrm{Ο}}y$ και $x\hat{\mathrm{Ο}}\delta=37^\circ$. Βρες το μέτρο της $x\hat{\mathrm{Ο}}y$.

\ex Σχεδίασε ένα τρίγωνο $\mathrm{ΑΒΓ}$ και φέρε τις διχοτόμους των γωνιών του. Τι παρατηρείς;

\ex Σωστό ή λάθος; α) Η διχοτόμος είναι ημιευθεία. \quad β) Η διχοτόμος χωρίζει τη γωνία σε δύο άνισα μέρη. \quad γ) Κάθε σημείο της διχοτόμου ισαπέχει από τις πλευρές της γωνίας. \quad δ) Η διχοτόμος γωνίας $80^\circ$ σχηματίζει με κάθε πλευρά γωνία $40^\circ$. \quad ε) Η διχοτόμος της ευθείας γωνίας είναι κάθετη στις πλευρές της.

\ex Σχεδίασε μια ευθεία $\varepsilon$ και ένα σημείο $\mathrm{Α}$ εκτός της. Φέρε με τον γνώμονα την κάθετη από το $\mathrm{Α}$ προς την $\varepsilon$ και μέτρησε την απόσταση του $\mathrm{Α}$ από την $\varepsilon$.

\ex Σχεδίασε γωνία $x\hat{\mathrm{Ο}}y=80^\circ$ και πάρε πάνω στη διχοτόμο της ένα σημείο $\mathrm{Α}$ με $\mathrm{ΟΑ}=5$ cm. Μέτρησε τις αποστάσεις του $\mathrm{Α}$ από τις $\mathrm{Ο}x$ και $\mathrm{Ο}y$.

\zoni{Το χρησιμοποιώ}{πίτσες, ρολόγια, δρόμοι}

\ex \textbf{Η πίτσα.} Μια φέτα πίτσα έχει γωνία κορυφής $50^\circ$. Ο Νίκος τη χωρίζει στα δύο με την διχοτόμο. Πόση γωνία έχει καθένα από τα δύο κομμάτια;

\ex \textbf{Το μονοπάτι.} Η γωνία της γωνιάς του φράχτη στον κήπο της θείας Μαρίας είναι $84^\circ$. Ένα μονοπάτι τη χωρίζει σε δύο ίσα παρτέρια. Με ποια γωνία ξεκινά το μονοπάτι από κάθε πλευρά του φράχτη;

\ex \textbf{Το ρολόι.} Στις $8{:}00$, ποιον αριθμό του ρολογιού δείχνει η διχοτόμος της (μικρότερης) γωνίας των δεικτών; Και στις $2{:}00$;

\ex \textbf{Ο ανεμιστήρας.} Ένας ανεμιστήρας της Αγγελικής έχει τρία πτερύγια που σχηματίζουν μεταξύ τους ίσες γωνίες. Πόση γωνία έχει η διχοτόμος δύο διαδοχικών πτερυγίων με καθένα από αυτά;

\ex \textbf{Το περίπτερο.} Δύο ευθείες οδοί τέμνονται σε γωνία $60^\circ$ και ένα περίπτερο πρέπει να ισαπέχει από τις δύο οδούς και να είναι $6$ cm από τη διασταύρωση στο σχέδιο. Σχεδίασε πού μπαίνει.

\zoni{Το σκέφτομαι}{γιατί και πώς}

\ex Γιατί κάθε γωνία έχει μία και μοναδική διχοτόμο; Τι γίνεται με τη διχοτόμο της μηδενικής γωνίας;

\ex Ένα σημείο στο εσωτερικό μιας γωνίας απέχει $3$ cm από τη μία πλευρά και $5$ cm από την άλλη. Ανήκει στη διχοτόμο; Πού να το μετακινήσεις ώστε να ανήκει;

\ex Δύο γωνίες είναι διαδοχικές με άθροισμα $180^\circ$. Η πρώτη είναι $50^\circ$. Σχεδίασε τις διχοτόμους και βρες τη γωνία που σχηματίζουν μεταξύ τους.

\ex Δες τον τρίτο τρόπο, του διαβήτη, για τη διχοτόμο. Σχεδίασε γωνία $70^\circ$ και κατασκεύασε τη διχοτόμο της χωρίς μοιρογνωμόνιο. Έλεγξέ την με μοιρογνωμόνιο.

\zoni{Για το σπίτι}{με χαρτί}

\ex Κόψε από χαρτί μια οξεία και μια αμβλεία γωνία και βρες τη διχοτόμο τους με δίπλωση. Μέτρησε με μοιρογνωμόνιο ότι οι δύο μισές είναι ίσες.

\ex Στο σπίτι σου βρες τρία πράγματα με γωνίες (π.χ. ένα βιβλίο ανοιχτό, ένα ψαλίδι) και πες πού θα περνούσε η διχοτόμος τους.
